% !TEX root = chapter-standalone.tex

\section{New monotone relations from old}
\label{sec:monrel-new-from-old}

In \cref{chap:poset-build} we constructed new \SY{preorders} from old ones: products, sums, and opposites.
We could construct products and sums of monotone functions; these constructions extend also to monotone relations.
The opposite of a preorder, in turn, is matched by the transpose of a relation.
Throughout this section, \posA, \posB, \posC, and \posD are \SY{preorders}.

\subsection{Products}

\begin{ctdefinition}[Product of monotone relations]
    \label{def:monrel-product}
    Let $\relA \colon \posA \profto \posB$ and $\relB \colon \posC \profto \posD$ be monotone relations.
    Their \emph{product} is the relation
    \begin{equation}
        \relA \Ctimes \relB \colon \posA \Ptimes \posC \profto \posB \Ptimes \posD
    \end{equation}
    between the product preorders (\cref{def:poset-product}) given by
    \begin{equation}
        \inrel{\tup{\posAel, \posCel}}{(\relA \Ctimes \relB)}{\tup{\posBel, \posDel}}
        \quad \text{if and only if} \quad
        \inrel{\posAel}{\relA}{\posBel}
        \ \text{ and } \
        \inrel{\posCel}{\relB}{\posDel}.
    \end{equation}
\end{ctdefinition}

In words: a pair is related to a pair if the first components are related by \relA and the second components are related by \relB.
The two components do not interact.

\begin{lemma}
    \label{lem:monrel-product}
    The product $\relA \Ctimes \relB$ of two monotone relations is a monotone relation.
\end{lemma}
\begin{proof}
    Recall that $\tup{\posAel', \posCel'} \posleqof{\posA \Ptimes \posC} \tup{\posAel, \posCel}$ means $\posAel' \posAleq \posAel$ and $\posCel' \posCleq \posCel$.
    \begin{enumerate}
        \item Closure downward in the source: suppose $\tup{\posAel', \posCel'} \posleqof{\posA \Ptimes \posC} \tup{\posAel, \posCel}$ and $\inrel{\tup{\posAel, \posCel}}{(\relA \Ctimes \relB)}{\tup{\posBel, \posDel}}$.
              Then $\posAel' \posAleq \posAel$ and $\inrel{\posAel}{\relA}{\posBel}$, so $\inrel{\posAel'}{\relA}{\posBel}$ since \relA is monotone; and $\posCel' \posCleq \posCel$ and $\inrel{\posCel}{\relB}{\posDel}$, so $\inrel{\posCel'}{\relB}{\posDel}$ since \relB is monotone.
              Hence $\inrel{\tup{\posAel', \posCel'}}{(\relA \Ctimes \relB)}{\tup{\posBel, \posDel}}$.
        \item Closure upward in the target: suppose $\inrel{\tup{\posAel, \posCel}}{(\relA \Ctimes \relB)}{\tup{\posBel, \posDel}}$ and $\tup{\posBel, \posDel} \posleqof{\posB \Ptimes \posD} \tup{\posBel', \posDel'}$.
              In the same way, closure upward in the target of \relA and of \relB gives $\inrel{\posAel}{\relA}{\posBel'}$ and $\inrel{\posCel}{\relB}{\posDel'}$, so $\inrel{\tup{\posAel, \posCel}}{(\relA \Ctimes \relB)}{\tup{\posBel', \posDel'}}$.
    \end{enumerate}
\end{proof}

\begin{example}
    The product of the identity monotone relations is the identity monotone relation of the product:
    \begin{equation}
        \posAleq \Ctimes \posCleq = \posleqof{\posA \Ptimes \posC}.
    \end{equation}
    This is just a restatement of \cref{def:poset-product}: a pair is below a pair exactly when each component is below the corresponding component.
\end{example}

The product of monotone relations is compatible with the product of monotone functions.
Given \SY{monotone functions} $\mapa \colon \posA \mto \posB$ and $\mapb \colon \posC \mto \posD$, their product is the function
\begin{equation}
    \defmapperiod{
        \mapa \Ctimes \mapb
    }{
        \posA \Ptimes \posC
    }{
        \mto
    }{
        \posB \Ptimes \posD
    }{
        \tup{\posAel, \posCel}
    }{
        \tup{\mapa(\posAel), \mapb(\posCel)}
    }
\end{equation}
It is monotone: if $\posAel' \posAleq \posAel$ and $\posCel' \posCleq \posCel$, then $\mapa(\posAel') \posBleq \mapa(\posAel)$ and $\mapb(\posCel') \posDleq \mapb(\posCel)$.

\begin{lemma}
    \label{lem:companion-product}
    For \SY{monotone functions} $\mapa \colon \posA \mto \posB$ and $\mapb \colon \posC \mto \posD$,
    \begin{equation}
        \comp{\mapa \Ctimes \mapb} = \comp{\mapa} \Ctimes \comp{\mapb}
     \end{equation}
     and
     \begin{equation}
        \conj{\mapa \Ctimes \mapb} = \conj{\mapa} \Ctimes \conj{\mapb}.
    \end{equation}
\end{lemma}
\begin{proof}
    For companions: $\inrel{\tup{\posAel, \posCel}}{\comp{\mapa \Ctimes \mapb}}{\tup{\posBel, \posDel}}$ means $\tup{\mapa(\posAel), \mapb(\posCel)} \posleqof{\posB \Ptimes \posD} \tup{\posBel, \posDel}$, that is, $\mapa(\posAel) \posBleq \posBel$ and $\mapb(\posCel) \posDleq \posDel$.
    This says exactly that $\inrel{\posAel}{\comp{\mapa}}{\posBel}$ and $\inrel{\posCel}{\comp{\mapb}}{\posDel}$.
    The argument for conjoints is the same, with the inequalities reversed.
\end{proof}

\subsection{Sums}

Recall that the elements of the sum $\posA \Pplus \posC$ (\cref{def:disjoint-union-of-posets}) are tagged elements $\tup{1, \posAel}$ with $\posAel \setin \posAset$ and $\tup{2, \posCel}$ with $\posCel \setin \posCset$, and that elements with different tags are never related.

\begin{ctdefinition}[Sum of monotone relations]
    \label{def:monrel-sum}
    Let $\relA \colon \posA \profto \posB$ and $\relB \colon \posC \profto \posD$ be monotone relations.
    Their \emph{sum} is the relation
    \begin{equation}
        \relA \Cplus \relB \colon \posA \Pplus \posC \profto \posB \Pplus \posD
    \end{equation}
    given by
    \begin{equation}
        \inrel{\tup{i, \elc}}{(\relA \Cplus \relB)}{\tup{j, \eld}}
        \quad \text{if and only if} \quad
        \begin{cases}
            i = j = 1 \text{ and } \inrel{\elc}{\relA}{\eld}, \text{ or} \\
            i = j = 2 \text{ and } \inrel{\elc}{\relB}{\eld}.
        \end{cases}
    \end{equation}
\end{ctdefinition}

In words: $\relA \Cplus \relB$ acts as \relA on the first summand and as \relB on the second summand, and never relates an element of one summand to an element of the other.

\begin{lemma}
    \label{lem:monrel-sum}
    The sum $\relA \Cplus \relB$ of two monotone relations is a monotone relation.
\end{lemma}
\begin{proof}
    \begin{enumerate}
        \item Closure downward in the source: suppose $\tup{i', \elc'} \posleqof{\posA \Pplus \posC} \tup{i, \elc}$ and $\inrel{\tup{i, \elc}}{(\relA \Cplus \relB)}{\tup{j, \eld}}$.
              By the definition of the sum preorder, $i' = i$; by the definition of $\relA \Cplus \relB$, $i = j$.
              If $i = 1$, then $\elc' \posAleq \elc$ and $\inrel{\elc}{\relA}{\eld}$, so $\inrel{\elc'}{\relA}{\eld}$ since \relA is monotone.
              If $i = 2$, the same argument works with $\posCleq$ and \relB.
              In either case, $\inrel{\tup{i', \elc'}}{(\relA \Cplus \relB)}{\tup{j, \eld}}$.
        \item Closure upward in the target: suppose $\inrel{\tup{i, \elc}}{(\relA \Cplus \relB)}{\tup{j, \eld}}$ and $\tup{j, \eld} \posleqof{\posB \Pplus \posD} \tup{j', \eld'}$.
              Then $i = j = j'$, and closure upward in the target of \relA (if $i = 1$) or of \relB (if $i = 2$) gives $\inrel{\tup{i, \elc}}{(\relA \Cplus \relB)}{\tup{j', \eld'}}$.
    \end{enumerate}
\end{proof}

\begin{example}
    As for products, the sum of the identity monotone relations is the identity monotone relation of the sum:
    \begin{equation}
        \posAleq \Cplus \posCleq = \posleqof{\posA \Pplus \posC}.
    \end{equation}
    This is a restatement of \cref{def:disjoint-union-of-posets}.
\end{example}

Again, the sum of monotone relations is compatible with the sum of monotone functions.
Given \SY{monotone functions} $\mapa \colon \posA \mto \posB$ and $\mapb \colon \posC \mto \posD$, their sum $\mapa \Cplus \mapb \colon \posA \Pplus \posC \mto \posB \Pplus \posD$ is defined by
\begin{equation}
    (\mapa \Cplus \mapb)(\tup{1, \posAel}) = \tup{1, \mapa(\posAel)}
    \qqand
    (\mapa \Cplus \mapb)(\tup{2, \posCel}) = \tup{2, \mapb(\posCel)}.
\end{equation}
It is monotone, because related elements of $\posA \Pplus \posC$ lie in the same summand, where $\mapa \Cplus \mapb$ acts as the monotone function $\mapa$ or $\mapb$.

\begin{lemma}
    \label{lem:companion-sum}
    For \SY{monotone functions} $\mapa \colon \posA \mto \posB$ and $\mapb \colon \posC \mto \posD$,
    \begin{equation}
        \comp{\mapa \Cplus \mapb} = \comp{\mapa} \Cplus \comp{\mapb}
     \end{equation}
     and
     \begin{equation}
        \conj{\mapa \Cplus \mapb} = \conj{\mapa} \Cplus \conj{\mapb}.
    \end{equation}
\end{lemma}
\begin{proof}
    For companions: consider $\tup{1, \posAel}$ in the first summand of the source.
    It is related by $\comp{\mapa \Cplus \mapb}$ to $\tup{j, \eld}$ if and only if $\tup{1, \mapa(\posAel)} \posleqof{\posB \Pplus \posD} \tup{j, \eld}$, that is, if and only if $j = 1$ and $\mapa(\posAel) \posBleq \eld$.
    This is exactly the condition for $\inrel{\tup{1, \posAel}}{(\comp{\mapa} \Cplus \comp{\mapb})}{\tup{j, \eld}}$.
    Elements of the second summand are treated in the same way, with $\mapb$ in place of $\mapa$, and the argument for conjoints is the same with the inequalities reversed.
\end{proof}

\subsection{Transpose}

Every relation $\relA \colon \setA \sto \setB$ has a transpose $\relA\reltransp \colon \setB \sto \setA$ (\cref{def:relation-transpose}).
The transpose of a monotone relation is again a monotone relation, provided that we also take the opposites (\cref{def:poset-opposite}) of the source and target.

\begin{lemma}
    \label{lem:monrel-transpose}
    Let $\relA \colon \posA \profto \posB$ be a monotone relation.
    Then its transpose is a monotone relation
    \begin{equation}
        \relA\reltransp \colon \posB\posop \profto \posAop.
    \end{equation}
\end{lemma}
\begin{proof}
    Recall that $\inrel{\posBel}{\relA\reltransp}{\posAel}$ means $\inrel{\posAel}{\relA}{\posBel}$.
    \begin{enumerate}
        \item Closure downward in the source $\posB\posop$: suppose $\posBel' \posleqopof{\posB} \posBel$, that is, $\posBel \posBleq \posBel'$, and $\inrel{\posBel}{\relA\reltransp}{\posAel}$, that is, $\inrel{\posAel}{\relA}{\posBel}$.
              Since \relA is closed upward in its target, $\inrel{\posAel}{\relA}{\posBel'}$, that is, $\inrel{\posBel'}{\relA\reltransp}{\posAel}$.
        \item Closure upward in the target $\posAop$: suppose $\inrel{\posBel}{\relA\reltransp}{\posAel}$ and $\posAel \posAleqop \posAel'$, that is, $\posAel' \posAleq \posAel$.
              Since \relA is closed downward in its source, $\inrel{\posAel'}{\relA}{\posBel}$, that is, $\inrel{\posBel}{\relA\reltransp}{\posAel'}$.
    \end{enumerate}
\end{proof}

Taking opposites is necessary here: transposing exchanges the roles of source and target, and hence exchanges the two conditions of \cref{def:monotone-relation}, one of which asks for closure \emph{downward} and the other for closure \emph{upward}.

\begin{lemma}
    \label{lem:monrel-transpose-properties}
    Let $\relA \colon \posA \profto \posB$ and $\relB \colon \posB \profto \posC$ be monotone relations.
    Then:
    \begin{enumerate}
        \item $(\relA\reltransp)\reltransp = \relA$;
        \item $(\relA \mthen \relB)\reltransp = \relB\reltransp \mthen \relA\reltransp$;
        \item $(\posAleq)\reltransp = \posAleqop$, so the transpose of the identity monotone relation on \posA is the identity monotone relation on $\posAop$.
    \end{enumerate}
\end{lemma}
\begin{proof}
    \begin{enumerate}
        \item Transposing twice exchanges the two components twice.
        \item $\inrel{\posCel}{(\relA \mthen \relB)\reltransp}{\posAel}$ means that there is $\posBel$ with $\inrel{\posAel}{\relA}{\posBel}$ and $\inrel{\posBel}{\relB}{\posCel}$, that is, with $\inrel{\posCel}{\relB\reltransp}{\posBel}$ and $\inrel{\posBel}{\relA\reltransp}{\posAel}$; and this means $\inrel{\posCel}{(\relB\reltransp \mthen \relA\reltransp)}{\posAel}$.
        \item $\inrel{\posAel}{(\posAleq)\reltransp}{\posAel'}$ means $\posAel' \posAleq \posAel$, which by \cref{def:poset-opposite} means $\posAel \posAleqop \posAel'$.
    \end{enumerate}
\end{proof}

The transpose connects the two ways of turning a monotone function into a monotone relation.
A \SY{monotone function} $\mapa \colon \posA \mto \posB$ is also a monotone function $\mapa\posop \colon \posAop \mto \posB\posop$ between the opposite preorders, with the same underlying function: if $\posAel \posAleqop \posAel'$, then $\posAel' \posAleq \posAel$, so $\mapa(\posAel') \posBleq \mapa(\posAel)$, that is, $\mapa(\posAel) \posleqopof{\posB} \mapa(\posAel')$.

\begin{lemma}
    \label{lem:conjoint-via-transpose}
    For every \SY{monotone function} $\mapa \colon \posA \mto \posB$,
    \begin{equation}
        \conj{\mapa} = \pars{\comp{\mapa\posop}}\reltransp,
    \end{equation}
    as monotone relations $\posB \profto \posA$.
\end{lemma}
\begin{proof}
    The companion $\comp{\mapa\posop} \colon \posAop \profto \posB\posop$ relates $\posAel$ to $\posBel$ if and only if $\mapa(\posAel) \posleqopof{\posB} \posBel$, that is, if and only if $\posBel \posBleq \mapa(\posAel)$.
    By \cref{lem:monrel-transpose}, its transpose is a monotone relation $(\posB\posop)\posop \profto (\posA\posop)\posop$, that is, $\posB \profto \posA$, and it relates $\posBel$ to $\posAel$ if and only if $\posBel \posBleq \mapa(\posAel)$.
    This is the definition of $\conj{\mapa}$.
\end{proof}

\begin{example}
    The transpose of the monotone relation $\relA$ of \cref{exa:monrel-coins} is the monotone relation $\relA\reltransp \colon \posB\posop \profto \posAop$ which relates a combination of coins to each amount it suffices to pay.
    In the opposite orders, ``fewer coins'' is now ``larger'', and ``smaller amount'' is now ``larger'': if some coins suffice to pay an amount, then more coins suffice to pay it, and the same coins suffice to pay any smaller amount.
\end{example}


